\documentclass[10pt,a4paper]{amsart}
\usepackage[margin=19mm]{geometry}
\usepackage{amsmath,amssymb,mathtools}
\usepackage[scr=rsfs,cal=euler]{mathalfa}
\usepackage{xfrac}
\usepackage[unicode,hidelinks,bookmarks=false]{hyperref}
\newcommand{\ie}{i.e.\ }
\newcommand{\cf}{cf.\ }
\newcommand{\resp}{respectively\ }
\newcommand{\Subsection}{\subsection}
\providecommand{\nobreakdash}{\nobreak\hbox{-}}
\setlength{\emergencystretch}{2em}
\multlinegap0pt
\usepackage{amssymb}
\usepackage{xcolor}
\usepackage{algorithm}\usepackage{algpseudocode}
\newtheorem{theorem}{Theorem}
\title{Norm Principles for ${\rm GO}(q)$ and ${\rm Spin}(q)$}
\author{Priyabrata Mandal}
\date{Source: arXiv:2608.28286v1 (CC BY 4.0)}
\begin{document}
\maketitle

\noindent\emph{Source and licence.} Norm Principles for ${\rm GO}(q)$ and ${\rm Spin}(q)$. Version arXiv:2608.28286v1 (CC BY 4.0), distributed by arXiv under the Creative Commons Attribution 4.0 International licence, http://creativecommons.org/licenses/by/4.0/, as recorded in the arXiv licence metadata for this exact version and shown on the abstract page https://arxiv.org/abs/2608.28286v1. Full licence notice: \texttt{licenses/arxiv-2608.28286-CC-BY-4.0.txt}.\par\smallskip

\noindent\textit{Editorial scope.} Counted unit: the article's theorem theorem (source0.tex lines 292--294). The article's complete original proof of it is delivered with it (source0.tex lines 295--301, one \texttt{proof} environment of 7 lines). No mathematics is written, repaired or re-derived: the statement and the proof are the article's own lines, in the article's own notation, and no retained line is substituted or re-flowed.  No further source-local context is needed: every cross-reference of every retained line resolves inside the two delivered spans.  Nothing was omitted from any retained span, so the package carries no omission record.  No retained line contains a citation, so the wrapper carries no bibliography and no bibliography entry.  No retained line contains a figure environment, an image-inclusion command or drawing code, so no image is delivered and the package declares no asset dependency.  Editorial formatting only, in addition: an AMS \texttt{amsart} preamble that restates the article's own declarations and macros used by the retained lines, namely the environments \texttt{theorem} on separate counters, together with the standard packages those lines need.  The retained environments print in this document as Theorem 1 in order of appearance, whereas the article prints the same items with the numbering of its own full document.  The complete native source of the article as distributed in the arXiv e-print for this exact version is bundled unchanged as \texttt{sources/208778/source0.tex}.
\begin{theorem}
    The elements of $W(F)$ are in one-one correspondence with the isometry classes of all anisotropic forms over $F$.
\end{theorem}

\begin{proof}
    Any element in $\widehat W(F)$ can be written as $[q_1]-[q_2]$, where $[q_1], [q_2]$ are isometry classes of regular quadratic forms of $q_1$ and $q_2$ respectively. We claim that any element in $W(F)$ is of the form $[q]$, where $q$ is a regular quadratic form over $F$. Indeed, $[q_1]-[q_2]= [q_1 \perp (-q_2)]-[q_2 \perp (-q_2)]$. Now, for any scalar $a \in F^\times$, 
    \begin{equation} \label{witt ring equation}
          \langle a \rangle \perp \langle -a \rangle \cong \langle a, -a \rangle \cong a \langle 1,-1 \rangle = a \mathbb H = 0 \in W(F)
    \end{equation}
    which implies $- \langle a \rangle= \langle -a \rangle \in W(F)$. Note that, $[q_2 \perp (-q_2)]$ is a hyperbolic form over $F$ and hence it becomes zero over $W(F)$. Therefore, $[q_1]-[q_2]= [q_1 \perp (-q_2)] \in W(F)$. If there is any common scalar in the diagonalization of $q_1$ and $q_2$, then by the above equation \ref{witt ring equation}, it denote the zero element in $W(F)$. In particular, every element in $W(F)$ is represented by a form $[q]$. By Witt decomposition theorem, every quadratic form $q$ can be written as $q=q_h \perp q_a$, where $q_h$ denotes the hyperbolic part of $q$ and $q_a$ denotes the anisotropic part of $q$. Therefore, $[q]$ and $[q_a]$ denotes the same element in $W(F)$. To prove the one-one correspondence, we need to show if $[q]$ and $[\phi]$ be two regular quadratic forms such that $[q]= [\phi ] \in W(F)$, then $q \cong \phi$. If $[q]= [\phi ] \in W(F)$, then $[q]= [\phi] \perp m \mathbb H \in \widehat W(F)$ for some integer $m \geq 0$. Hence, $q \cong \phi \perp m \mathbb H$ as quadratic forms. Since, $q$ is anisotropic, it cannot contain a hyperbolic subpart in it, and so $m=0$. Thus, $q \cong \phi$. This completes the proof.
\end{proof}

\end{document}
